\documentclass[10pt]{article}

\usepackage[utf8]{inputenc}

\usepackage[T1]{fontenc}

\usepackage{amsmath}

\usepackage{amsfonts}

\usepackage{amssymb}

\usepackage[version=4]{mhchem}

\usepackage{stmaryrd}

\usepackage{mathrsfs}

\usepackage{bbold}

\usepackage{graphicx}

\usepackage[export]{adjustbox}

\graphicspath{ {./images/} }



\begin{document}

\begin{enumerate}

  \setcounter{enumi}{1}

  \item Т е о р мы о не подвижных то чк а х. a) Пусть $E$ - отделимое л. в. П., $K$ - его непустое выпуклое компактное подмножество и $f$ - отображение множества $K$ в множество его непустых выпуклых замкнутых подмножеств. Пусть для каждого $x \in K$ и каждой окрестности $U$ множества $f(x)$ существует такая окрестность $\mathscr{P}$ точки $x$, что $f(\mathscr{V} \cap K) \subset \mathcal{U}$ (это свойство отображения $f$ наз. его по л у н е п р е р ы вн о с т ь ю с в е p $\mathrm{x} \mathrm{y}$ ). Тогда существует точка $z \in K$ такая, что $z \in f(z)$ - "неподвижная точка» отображения $f$ (т е о р м а Ф а н я-обобшение те о ре мы III а у д е р а - Тихонов а). б) Пусть $E$ - отделимое Т. в. П., $K$ - его непустое компактное выпуклое подмножество и $\Gamma$ - некоторое множество попарно коммутирующих непрерывных отображений множества $K$ в $K$, обладающее следующим свойством: если $x$, $z \in K, \quad \alpha, \quad \beta \in \mathbb{R}, \quad \alpha, \beta>0, \quad \alpha+\beta=1, \quad$ то $g(\alpha x+\beta z)=$ $=\alpha g(x)+\beta g(z)$. Тогда существует точка $z_{0} \in K$ такая, что $g\left(z_{0}\right)=z_{0}$ для всех $g \in \Gamma \quad$ (т е о р е м а M а р к ов а - К а к у т ан и).



  \item Важное значение в теории ЛВП имеют также Хана - Банаха теорема и Банаха - IIIтейнхауза теорема.



\end{enumerate}



Ряд интересных результатов получен в теории мер, принимающих значения в л. в. П. и (особенно) в связанной с теорией случайных процессов теории числовых цилиндрич. мер, определенных на л. в. п.



Возник и продолжает развиваться математич. анализ (в широком смысле) на Т. в. П.- т. н. б е с к о н е чн о м е р ны й а н а ли з. Являясь продолжением классич. анализа, он в то же время качественно отличается от него как постановками задач и результатами, так и методами. Бесконечномерный анализ включает теорию диффференцируемых отображений Т. в. п. и диффференцируемых мер на Т. в. п.; теорию обобщенных функций и мер (распределений) на Т. в. п.; теорию дифференциальных уравнений - как относительно функций вещественного аргумента, принимающих значения в Т. в. п., так и относительно числовых функций и мер (возможно, обобщенных), определенных на Т. в. п. На языке бесконечномерного анализа весьма естественно формулируются основные задачи физики бесконечномерных систем - квантовой теории поля, статистич. механики, гидродинамики, - а также нек-рые математич. задачи, первоначально возникшие вне бесконечномерного анализа.



Лит.: [1] Б у р б а к и Н., Топологические векторные пространства, пер. с франц., М., 1959; [2] Р о б е р т с о н А. Р о б е р т с он В., Топологические векторные пространства, пер. с англ., М., 1967; [3] Ш е ф е р Х., Топологические векторные пространства, пер. с англ., М., 1971; [4] Э д в а р д с Р.-Э. Функциональный, анализ, пер. с англ., М., 1969; [5] П и ч А., Ядерные локально выпуклые пространства, пер. с нем., М., 1967; [6] е го ж е, Операторные идеалы, пер. с англ., М., 1982; [7] W il d e M. d e, Closed graph theorems and webbed spaces, L. $1978 ;$ [8] S c h w a r t z L., Théorie des distributions, [2 éd.], $\mathrm{P}$. 1978; [8] S c h w a r t z L., Théorie des distributions, [2 éd.], P.,



\includegraphics[max width=\textwidth, center]{2023_07_09_26884c2e1ee13e11741fg-1}

"Функц. анализ и его прил.", 1977, т. 11, в. 1, с. $91-92$; [10] С м о л я н о в О. Г , “Изв. АН СССР. Сep. матем.», 1971, т. 35, № 3, с. $682-96$; [11] е г о ж е, Анализ на топологических линейных пространствах и его приложения, М., 1979.



ТОПОЛОГИ ЧЕСКОЕ КОЛЬЦО - кольцо $\stackrel{R}{R,}$ я вляющееся топологич. пространством, шричем требуется, ттобы отображения



\[

(x, y) \rightarrow x-y \text { и }(x, y) \rightarrow x y

\]



были непрерывны. Т. к. $R$ наз. о т д е л и м ы м, если оно отделимо как топологич. пространство. В этом случае пространство $R$ хаусдорфово. Любое подкольцо $M$ Т.к. $R$, а также факторкольцо $R / J$ по идеалу $J$ являются Т. к. Если $R$ отделимо и идеал $J$ замкнут, то $R / J$ - отделимое Т. к. Замыкание $\bar{M}$ подкольца $M$ в $R$ также является Т. к. Прямое произведение топологич. колеці - Т. к. Го м о м о ф и з топологич. колец - это гомоморфизм колец, являющийся непрерывным отображением. Если $f: R_{1} \rightarrow R_{2}$ - такой гомоморфизм, причем $f$ - эпиморфизм и открытое отображение, то $R_{2}$ как Т. к. изоморфно кольцу $R_{1} / \operatorname{Ker} f$. Примеры Т. к. доставляют банаховы алгебры. Важный тип Т. к. определяется тем условием, что в качестве фундаментальной системы окрестностей нуля можно выбрать некоторое множество идеалов. Например, с любым идеалом $\mathfrak{m}$ коммутативного кольца $R$ связана $\mathfrak{m}$-адическая топология, в $\kappa$-рой множества $\mathfrak{m}^{n}$ для всех натуральных $n$ образуют фундаментальную систему окрестностей нуля. Эта топология отделима, если выполнено условие



\[

\bigcap_{n} \mathfrak{m}^{n}=0

\]



Для Т. к. $R$ определено его пополнение $\hat{R}$, являющееся полным Т. к., причем отделимое кольцо $R$ вкладывается в $\hat{R}$, к-рое тоже отделимо, как всюду плотное подмножество. Аддитивная группа кольца $\widehat{R}$ совпадает с пополнением аддитивной грушпы кольца $R$ как абелевой топологит. группы.



Лum.: [1] Б у р б а к и Н., Общая топология. Топологические группы. Числа и связанные с ними группы и пространства, пер. с франц., М., 1969; [2] е г о ж е, Коммутативная алгебра, пер. с франц., М., 1971; [3] П о н т р я г и н Ј. С., Непрерывные группы, 3 изд., М., $1973 ;[4]$ В а н д е р В а р д е н Б. Ј.,



ТОПОЛОГИЧЕСКОЕ ПОЛЕ - топологическое кольцо $K$, являющееся полем, причем дополнительно требуется, чтобы отображение $a \rightarrow a^{-1}$ было непрерывно на $K \backslash\{0\}$. Любое подполе $P$ Т. п. $K$ и замыкание $\overline{\boldsymbol{P}}$ поля $P$ в $K$ снова являются Т. П.



Связные локально компактные Т. п. исчерпываются полями $\mathbb{R}$ и $\mathbb{C}$ (см. Локально компактное тело). Каждое нормированное поле является Т. п. относительно топологии, порождаемой нормой (см. Нормирование, А бсолютное значение). Если существуют два вещественных нормирования $u$ и $v$ поля $P$, каждое из к-рых превращает $P$ в полное Т. п., и топологии $\tau_{u}$ и $\tau_{v}$, порождаемые $u$ и $v$, различны, то поле $\boldsymbol{P}$ алгебраически замкнуто. Поле $\mathbb{C}$ - единственное вещественно нормированное расширение поля $\mathbb{R}$.



На каждом поле бесконечной мощности $\tau$ существует ровно $2^{2^{\tau}}$ различных топологий, превращающих его в Т. п. Топология Т. п. либо антидискретна, либо вполне регулярна. Построены Т. п., топология к-рого не нормальна, и Т. п., топология к-рого нормальна, но не наследственно нормальна. Т. П. либо связно, либо вполне несвязно. Существует связное Т. П. любой конеqной характеристики. Неизвестно (1983), можно ли каждое Т. П. вложить в связное Т. П. в качестве подполя. В отличие от топологич. колец и линейных топологич. пространств, не каждое вполне регулярное топологич. пространство можно вложить в качестве подпространства в Т. п. Так, напр., псевдокомпактное (в частности, бикомпактное) подпространство Т.п. всегда метризуемо. Однако каждое вполне регулярное пространство, допускающее взаимно однозначное непрерывное отображение на метрич. пространство, вкладывается в нек-рое Т. П. в качестве подпространства. Если в Т. ІІ. $P$ есть хоть одно счетное незамкнутое множество, то существует более слабая метризуемая топология на поле $P$, превращающая его в Т. Ш.



Для Т. п. $K$ определено его пополнение $\widetilde{K}-$ полное топологич. кольцо, в к-рое $K$ вкладывается как всюду плотное подполе. Кольцо $\tilde{K}$ может иметь делители нуля. Однако пополнение всякого вещественно нормированного Т. П. есть вещественно нормированное Т. п. Лит.: [1] П он т р я г и н Л. С., Непрерывные группы, 3 изд., M., 1973; [2] W i e s \& a w W., Topological fields, Wroc1983 , т. 271 , № 6, с. $1332-36$. Д. Б., "ДолЈ. Б. Шахматов.



$\triangle 13$ Математическая әнц., т. 5





\end{document}